Gemini models are trained on a dataset that is both multimodal and multilingual. Our pre-training dataset uses data from web documents, books, and code, and includes image, audio, and video data.

We use the SentencePiece tokenizer~\cite{sentencepiece} and find that training the tokenizer on a large sample of the entire training corpus improves the inferred vocabulary and subsequently improves model performance. For example, we find Gemini models can efficiently tokenize non-Latin scripts which can, in turn, benefit model quality as well as training and inference speed.

The number of tokens used to train the largest models were determined following the approach in \citet{chinchilla}. The smaller models are trained for significantly more tokens to improve performance for a given inference budget, similar to the approach advocated in~\citet{llama}. 

We apply quality filters to all datasets, using both heuristic rules and model-based classifiers. We also perform safety filtering to remove harmful content based on our policies. To maintain the integrity of evaluations, we search for and remove any evaluation data that may have been in our training corpus before using data for training. The final data mixtures and weights were determined through ablations on smaller models. We stage training to alter the mixture composition during training – increasing the weight of domain-relevant data towards the end of training. We find that data quality is an important factor for highly-performing models, and believe that many interesting questions remain around finding the optimal dataset distribution for pre-training.